% --------------------------------- Abstract ---------------------------------
\thispagestyle{plain}
\frontheading{Abstract}
\addcontentsline{toc}{chapter}{Abstract}

{\fontsize{12}{20}\selectfont
\setlength{\parindent}{0.5in}

India recognises twenty-two scheduled languages, yet nearly every machine
translation service that serves them well is a cloud service. That design
excludes precisely the users who need translation most: a health worker in a
village with no signal, a soldier on a border patrol, a citizen at a gram
panchayat counter. Cloud translation also means every sentence a user types
leaves their device. The obvious alternative, running a large translation model
locally, fails on arithmetic: state-of-the-art Indic systems such as IndicTrans2
need GPU-class memory and hold billions of parameters, which no commodity phone
or ARM board can host.

\setu{} (Sanskrit \dev{सेतु}, ``bridge'') closes that gap by distilling a strong
teacher into compact students that run entirely on the device. The project began
from the hypothesis that \emph{preference distillation} --- ranking teacher
hypotheses by chrF and training the student with Direct Preference Optimization
(DPO) --- would beat standard sequence-level knowledge distillation for tiny
translation models. A controlled head-to-head at matched student size refuted
that hypothesis and produced the project's central empirical finding: at 100{,}000
training sentences, sequence-level knowledge distillation (\seqkd{}) on the
teacher's own one-best output reached \textbf{17.09 BLEU} against
\textbf{11.10 BLEU} for reference-trained DPO, an improvement of 6.0 BLEU
(roughly 54\,\%). The bottleneck was never the optimisation objective; it was
noise in the mined human references. That result redirected the entire system
towards \seqkd{}, and the DPO machinery was retained as a measured, reported
negative result rather than quietly discarded.

The delivered system trains one 52.6-million-parameter Marian student per
translation direction on approximately 497{,}000 teacher-distilled sentence
pairs, then quantises it to INT4 ONNX at \textbf{103.9\,MB} --- a fourfold
compression from 411\,MB FP32. Twenty-two directions covering eleven scheduled
languages in both directions were trained on an NVIDIA A40 vGPU. Measured
against the teacher on a held-out Samanantar development slice,
\textbf{14 of 22 directions meet the $\geq 0.80$ student/teacher BLEU ratio
target}, with a mean ratio of \textbf{0.838} and a median of \textbf{0.855};
the strongest direction, Marathi\arrowto{}English, reaches \textbf{0.955}. All
three non-quality targets pass without exception: every artifact is 103.9\,MB
(target $\leq 200$\,MB), mean latency is 105--202\,ms per sentence on commodity
CPU (target $<500$\,ms), and inference completes with all network sockets
disabled. Because no Indic-to-Indic student exists, an English-pivot router
chains two students, so all 11 languages translate in every direction ---
22 direct models serve 132 usable directions. Two training collapses
(Telugu\arrowto{}English and Malayalam\arrowto{}English, both degenerating to a
fixed output string) were diagnosed as learning-rate instability at
$5\times10^{-4}$ and fully recovered at $3\times10^{-4}$, lifting ratios from
0.036 to 0.817 and from 0.013 to 0.898 respectively. The system ships as a
FastAPI service, an offline progressive web application, a Next.js web client
and a command-line tool, all sharing a single inference engine, and is validated
by an automated suite of 85 tests.

\vspace{10pt}

\noindent\textbf{Keywords:} Neural Machine Translation; Sequence-Level Knowledge
Distillation; Direct Preference Optimization; Small Language Models; Indic
Languages; On-Device Inference; Model Quantisation; ONNX; English Pivot.

}
\clearpage
